% GENERATED FILE. Edit canonical lessons, then run npm run generate:books.
% canonical-source-hash: fnv1a64:576d7e8fd10e3c0b
% canonical-lessons: ZH-C195-cunqian, ZH-C195-quqian, ZH-C195-xinyongka, ZH-C195-mima, ZH-C195-zhanghu

\chapter{At the Bank}
\label{ch:zh-at-the-bank}
\hlchaptermodality{195}

\begin{chapteropening}
\textbf{By the end of this chapter:} \emph{I can put money in and take it out, and talk about cards, PINs and bank accounts.}

Complete the last lesson of chapter 195: I can put money in and take it out, and talk about cards, PINs and bank accounts.
\end{chapteropening}

\section[cúnqián]{cúnqián — to put money by; to put money in the bank}

\label{lesson:ZH-C195-cunqian}



\begin{quote}
Before the new one: say the Mandarin for public square, then the Mandarin for traffic lights.
\end{quote}

\subsection*{You'll want to know: cúnqián}
\emph{cúnqián} — \textquotedblleft{}to put money by; to put money in the bank\textquotedblright{}.

\begin{quote}
\emph{wǒ měi ge yuè dōu cúnqián} — I put money aside every month.
\end{quote}

Its characters come later; for now it is a word for the ear.

\begin{cousinweb}
\emph{cúnqián} is built from \emph{cún} 'to keep' and \zh{钱} 'money', both inherited from Middle Chinese, \emph{dzwon} and \emph{dzjen} (Guangyun, in Baxter's transcription). The dictionary glosses the first 'to be present' and 'to care for'; putting money away is not given there.
\end{cousinweb}

\begin{grammarlens}[title={What you've built}]
One more word for reading notices and messages.
\end{grammarlens}

\subsection*{Your turn}
\begin{itemize}
\raggedright
  \item \emph{Say it:} \emph{cúnqián}
  \item \emph{Say it:} \emph{cúnqián}, once more
  \item \emph{Say it:} \emph{wǒ měi ge yuè dōu cúnqián} — I put money aside every month.
  \item \emph{Recall:} say the Mandarin for public square, then the Mandarin for traffic lights, then say \emph{cúnqián} again
\end{itemize}

\begin{tcolorbox}[breakable,colback=teal!4,colframe=teal!35!black,title={Before you move on}]
What does cúnqián mean? (\textquotedblleft{}To put money by; to put money in the bank\textquotedblright{}.) Say it once more.
\end{tcolorbox}
\section[qǔqián]{qǔqián — to take out money}

\label{lesson:ZH-C195-quqian}



\begin{quote}
Before the new one: say the Mandarin for traffic lights, then the Mandarin for to put money by.
\end{quote}

\subsection*{You'll want to know: qǔqián}
\emph{qǔqián} — \textquotedblleft{}to take out money\textquotedblright{}.

\begin{quote}
\emph{wǒ yào qù yínháng qǔqián} — I'm going to the bank to get some money out.
\end{quote}

Its characters come later; for now it is a word for the ear.

\begin{cousinweb}
\emph{qǔqián} is built from \emph{qǔ} 'to take' and \zh{钱} 'money', both inherited from Middle Chinese, \emph{tshjuX} and \emph{dzjen} (Guangyun, in Baxter's transcription), which glosses the first 'to gather in' and 'to receive'.
\end{cousinweb}

\begin{grammarlens}[title={What you've built}]
One more word for reading notices and messages.
\end{grammarlens}

\subsection*{Your turn}
\begin{itemize}
\raggedright
  \item \emph{Say it:} \emph{qǔqián}
  \item \emph{Say it:} \emph{qǔqián}, once more
  \item \emph{Say it:} \emph{wǒ yào qù yínháng qǔqián} — I'm going to the bank to get some money out.
  \item \emph{Recall:} say the Mandarin for traffic lights, then the Mandarin for to put money by, then say \emph{qǔqián} again
\end{itemize}

\begin{tcolorbox}[breakable,colback=teal!4,colframe=teal!35!black,title={Before you move on}]
What does qǔqián mean? (\textquotedblleft{}To take out money\textquotedblright{}.) Say it once more.
\end{tcolorbox}
\section[xìnyòngkǎ]{xìnyòngkǎ — credit card}

\label{lesson:ZH-C195-xinyongka}



\begin{quote}
Before the new one: say the Mandarin for to put money by, then the Mandarin for to take out money.
\end{quote}

\subsection*{You'll want to know: xìnyòngkǎ}
\emph{xìnyòngkǎ} — \textquotedblleft{}credit card\textquotedblright{}.

\begin{quote}
\emph{zhèr kěyǐ yòng xìnyòngkǎ ma} — Can I pay by credit card here?
\end{quote}

Its characters come later; for now it is a word for the ear.

\begin{cousinweb}
\emph{xìnyòngkǎ} is built from \emph{xìnyòng} 'trust, credit' and \emph{kǎ} 'a card'. The first two parts are inherited from Middle Chinese, \emph{sinH} and \emph{yowngH} (Guangyun, in Baxter's transcription), glossed 'true' and 'to use'; \emph{kǎ} is usually traced to English \emph{card}.
\end{cousinweb}

\begin{grammarlens}[title={What you've built}]
One more word for reading notices and messages.
\end{grammarlens}

\subsection*{Your turn}
\begin{itemize}
\raggedright
  \item \emph{Say it:} \emph{xìnyòngkǎ}
  \item \emph{Say it:} \emph{xìnyòngkǎ}, once more
  \item \emph{Say it:} \emph{zhèr kěyǐ yòng xìnyòngkǎ ma} — Can I pay by credit card here?
  \item \emph{Recall:} say the Mandarin for to put money by, then the Mandarin for to take out money, then say \emph{xìnyòngkǎ} again
\end{itemize}

\begin{tcolorbox}[breakable,colback=teal!4,colframe=teal!35!black,title={Before you move on}]
What does xìnyòngkǎ mean? (\textquotedblleft{}Credit card\textquotedblright{}.) Say it once more.
\end{tcolorbox}
\section[mìmǎ]{mìmǎ — password; PIN}

\label{lesson:ZH-C195-mima}



\begin{quote}
Before the new one: say the Mandarin for to take out money, then the Mandarin for credit card.
\end{quote}

\subsection*{You'll want to know: mìmǎ}
\emph{mìmǎ} — \textquotedblleft{}password; PIN\textquotedblright{}.

\begin{quote}
\emph{wǒ wàngle mìmǎ} — I've forgotten my PIN.
\end{quote}

Its characters come later; for now it is a word for the ear.

\begin{cousinweb}
\emph{mìmǎ} is built from \emph{mì} 'secret' and \emph{mǎ} 'a numeral', both inherited from Middle Chinese, \emph{mit} and \emph{maeX} (Guangyun, in Baxter's transcription). The dictionary glosses the first 'a ridge' and 'quiet', and the second as a stone like jade; how they came to mean secret and a number sign is not settled.
\end{cousinweb}

\begin{grammarlens}[title={What you've built}]
One more word for reading notices and messages.
\end{grammarlens}

\subsection*{Your turn}
\begin{itemize}
\raggedright
  \item \emph{Say it:} \emph{mìmǎ}
  \item \emph{Say it:} \emph{mìmǎ}, once more
  \item \emph{Say it:} \emph{wǒ wàngle mìmǎ} — I've forgotten my PIN.
  \item \emph{Recall:} say the Mandarin for to take out money, then the Mandarin for credit card, then say \emph{mìmǎ} again
\end{itemize}

\begin{tcolorbox}[breakable,colback=teal!4,colframe=teal!35!black,title={Before you move on}]
What does mìmǎ mean? (\textquotedblleft{}Password; PIN\textquotedblright{}.) Say it once more.
\end{tcolorbox}
\section[zhànghù]{zhànghù — bank account}

\label{lesson:ZH-C195-zhanghu}



\begin{quote}
Before the new one: say the Mandarin for credit card, then the Mandarin for password.
\end{quote}

\subsection*{You'll want to know: zhànghù}
\emph{zhànghù} — \textquotedblleft{}bank account\textquotedblright{}.

\begin{quote}
\emph{nǐ yǒu yínháng zhànghù ma} — Do you have a bank account?
\end{quote}

Its characters come later; for now it is a word for the ear.

\begin{cousinweb}
\emph{zhànghù} is built from \emph{zhàng} 'accounts' and \emph{hù} 'a household'. The second is inherited from Middle Chinese \emph{huX} (Guangyun, in Baxter's transcription), glossed as a door, half of a gate; no dictionary we use records an older form of the first, so its history is not settled.
\end{cousinweb}

\begin{grammarlens}[title={What you've built}]
One more word for reading notices and messages.
\end{grammarlens}

\subsection*{Your turn}
\begin{itemize}
\raggedright
  \item \emph{Say it:} \emph{zhànghù}
  \item \emph{Say it:} \emph{zhànghù}, once more
  \item \emph{Say it:} \emph{nǐ yǒu yínháng zhànghù ma} — Do you have a bank account?
  \item \emph{Recall:} say the Mandarin for credit card, then the Mandarin for password, then say \emph{zhànghù} again
\end{itemize}

\begin{tcolorbox}[breakable,colback=teal!4,colframe=teal!35!black,title={Before you move on}]
What does zhànghù mean? (\textquotedblleft{}Bank account\textquotedblright{}.) Say it once more.
\end{tcolorbox}
